\documentclass[../report.tex]{subfiles}

\begin{document}

\subsection{Windows Scanning and Buffer Overflow}
\subsubsection{Conditions}
The target machine (\textit{IP: 192.168.61.45}) was a Windows 7 machine with a 32-bit architecture running a vulnerable server process (VulnServer.exe) on port 5555. The machine was accessible via remote desktop protocol, enabling participants to observe processes via a debugger.

\subsubsection{Goals}
It was not part of the mission objective to stay unnoticed by any potential intrusion detection systems or system administrators, therefore not necessitating the participants to obscure their traces.
\paragraph{Active Scan}
The following information was to be gathered about the target machine:
\begin{itemize}
\begin{multicols}{2}
\item{Operating System}
\item{Open Ports}
\item{Service information}
\item{Vulnerabilities}
\end{multicols}
\end{itemize}

\paragraph{Exploiting Buffer Overflow Vulnerability}
The process \enquote{VulnServer} has a Buffer Overflow vulnerability.
The goal was to exploit that vulnerability to create a reverse shell on the target system.
The following qualities must be met:
\vspace{-3mm}
\begin{enumerate}
\item{The input should override the Extended Instruction Pointer with a suitable memory address.}
\item{The input contains code for a reverse shell that runs on the laboratory machine \enquote{Anwendungssicherheit} / Kali Linux on port 443}
\item{The input should not contain any bad characters that could hinder the attack.}
\end{enumerate}

\end{document}